\section{The price of resubmission-proof admission}\label{sec:theory}
Proposition~\ref{prop:lineage} shows that spending $\alpha$ over a lineage is sufficient. This section asks what that protection costs an honest author. Repeated submission until success is a known threat to evaluation: it motivates reliable leaderboards and the reusable holdout \citep{blum2015ladder,dwork2015reusable}, it has been documented in public LLM leaderboards \citep{singh2025leaderboard}, and strategic submitters are studied in principal-agent testing \citep{bates2022principal}. Two ingredients are known. Some form of spending is necessary (Proposition~\ref{thm:lineage}, an elementary observation alongside classical alpha-spending, e.g.\ \citealp{tian2021online}). One exam at level $a$ needs about $\log(1/a)/K_{\inf}$ variants, where $K_{\inf}$ is the information each variant carries against the closest useless skill (Proposition~\ref{thm:onexam}, from \citealp{agrawal2025stopping,honda2010asymptotically,larsson2025numeraire}). Our contribution is Theorem~\ref{thm:price}: combining the two, the extra cost of resubmission-proofness is the entropy of the honest attempt-count distribution divided by $K_{\inf}$, no lineage-valid gate whose exams are fixed in advance can charge less, and a mismatched spending schedule adds a Kullback-Leibler term in both the upper and the lower bound. Choosing the schedule is then a coding problem: Proposition~\ref{prop:minimax} gives the minimax schedule over a class of attempt-count laws, Theorem~\ref{thm:learned} a schedule learned from earlier lineages, and Theorem~\ref{thm:adaptive} shows that the price holds, to first order, for a gate that does not know the skill's effect. Proofs are in Appendix~\ref{app:theory}.

\paragraph{Scope, in brief.} Three points frame what follows. (i) The lower bounds concern gates whose exam $k$ is fixed before the lineage starts; we claim no lower bound for gates that carry evidence or levels from one version to the next, or across lineages. (ii) The cost $C$ counts the variants of the exam that admits the helpful version; the useless versions are stopped by the futility bet at its own fixed level, so the schedule controls only the entropy term of the total. (iii) The deployed level $\min(\alpha_t,\alpha\gamma^{\text{lin}}_k)$ composes e-LOND across the library with spending within a lineage (Proposition~\ref{prop:lineage}); the theorems price the lineage term. Each point is made precise in ``What the theorem covers'' below.

\paragraph{Setting.} Write $X_i=(1+D_i)/2\in[0,1]$ and $m_0=(1+\delta)/2$, as in Section~\ref{sec:gate}. An author submits versions $k=1,2,\dots$ of a skill, and version $k$ may depend on the whole history $\mathcal H_{k-1}$: earlier transcripts, decisions and the author's own randomness. Exam $k$ draws its variants after its commitment, so they are fresh given $\mathcal H_{k-1}$.
\begin{definition}[Conditional null across a lineage]\label{def:condnull}
A lineage is null if, for every $k$, conditional on $\mathcal H_{k-1}$, on version $k$ and on the registry's history $\mathcal G_{k-1}$ when exam $k$ starts (which contains $\mathcal H_{k-1}$ and fixes the deployed level), the scores of exam $k$ satisfy $\mathbb E[X_i\mid\mathcal F_{i-1},\mathcal G_{k-1}]\le m_0$.
\end{definition}
This is the assumption that Propositions~\ref{prop:valid} and~\ref{prop:lineage} use. It fails if an author can see exam $k$'s hidden variants before submitting version $k$, which is what the commitment before the beacon prevents. For a law $P$ on $[0,1]$ with mean above $m_0$, let
\begin{align*}
K_{\inf}(P,m_0)&=\inf\{\mathrm{KL}(P\,\|\,Q):\ Q \text{ a law on }[0,1],\ \mathbb E_QX\le m_0\},\\
g_P(\lambda)&=\mathbb E_P\log(1+\lambda(X-m_0)).
\end{align*}

\begin{proposition}[Lineage control: sufficient and tight]\label{thm:lineage}
\emph{(a)} If exam $k$ runs at a level $a_k$ fixed in advance and is valid under Definition~\ref{def:condnull}, then every adaptive author with a null lineage is ever admitted with probability at most $1-\prod_k(1-a_k)\le\sum_ka_k$.
\emph{(b)} Consider a gate whose exams are fixed in advance: the level, bet and stopping rule of exam $k$ do not depend on earlier exams' outcomes, and each exam draws its variants independently of earlier exams. Fix a null law $Q$ and let $\sigma_k(Q)$ be the probability that exam $k$ admits a version whose scores are i.i.d.\ $Q$. The author who resubmits such a version until it is admitted succeeds with probability $1-\prod_k(1-\sigma_k(Q))$. A gate that controls lineage error at $\alpha$ therefore needs $\sum_k\sigma_k(Q)\le\bar\alpha:=\log\frac{1}{1-\alpha}$ for every null $Q$. In particular, any gate whose exams all have size at least $\sigma>0$ under some null law admits a resubmitted useless skill with probability $1-(1-\sigma)^K\to1$.
\end{proposition}
Part (a) is the classical alpha-spending argument \citep{tian2021online} stated under Definition~\ref{def:condnull}. Part (b) is elementary; we state it because it shows that per-exam sizes must be summable, so no gate that leaves each exam at a fixed level can be resubmission-proof. The \NResubPerExam\% failure after \PResub{} resubmissions in Section~\ref{sec:res-library} is this effect: it equals a constant per-exam size of \NResubImpliedSize\%.

\begin{proposition}[Cost of one exam]\label{thm:onexam}
Let $X_1,X_2,\dots$ be i.i.d.\ $P$ with $\mathbb E_PX>m_0$.
\emph{(a)} Any exam that admits with probability at most $s$ under every i.i.d.\ null law and with probability at least $1-\beta$ under $P$ has $\mathbb E_P[\tau]\ge\{(1-\beta)\log(1/s)-\log2\}/K_{\inf}(P,m_0)$.
\emph{(b)} For any $\lambda\in[0,1/m_0]$ with $g_P(\lambda)>0$, the exam that stops when $\prod_{i\le n}(1+\lambda(X_i-m_0))\ge1/a$ is valid at level $a$, has power one, and has $\mathbb E_P[\tau]\le\{\log(1/a)+\log(1/m_0)\}/g_P(\lambda)$.
\emph{(c)} $\sup_{0\le\lambda<1/m_0}g_P(\lambda)=K_{\inf}(P,m_0)$. The supremum is attained by a least-favourable null law $Q^*$ with $Q^*(x)\propto P(x)/(1+\lambda^*(x-m_0))$ on the support of $P$, plus any leftover mass at $0$.
\end{proposition}
None of this is new in substance. For power-one tests, \citet[Theorem~3.1]{agrawal2025stopping} prove the lower bound $\log(1/\alpha)/\mathrm{KL}_{\inf}$ for general composite nulls, with asymptotically matching upper bounds, and \citet{waudbysmith2025instance} prove matching asymptotic bounds for portfolio e-processes, bounded means included. Part (c) is the bounded-support duality of \citet{honda2010asymptotically} \citep[see also][Proposition~1]{honda2015nonasymptotic}, mirrored for a null mean bounded above, and $Q^*$ is the reverse information projection of \citet{larsson2025numeraire}. Part (a) uses the change-of-measure inequality of \citet[Lemma~1]{kaufmann2016complexity} \citep[see also][]{garivier2019explore}. We state the bounded case with explicit finite-level constants, allowing power $1-\beta<1$ as Theorem~\ref{thm:price} needs, and give a short proof that exhibits $Q^*$. Together the parts say that the cheapest valid exam at level $a$ needs $\log(1/a)/K_{\inf}$ variants, up to an additive $\log(2/m_0)/K_{\inf}$, and that the fixed Kelly bet $\lambda^*=\arg\max_\lambda g_P(\lambda)$ attains this \citep[compare][]{grunwald2024safe,shekhar2024nonparametric}.

\begin{theorem}[Price of resubmission-proofness]\label{thm:price}
An honest author's versions $1,\dots,J-1$ are useless and version $J$ has i.i.d.\ scores from $P$, where the attempt index $J$ has law $\pi$ and each event $\{J=k\}$ is fixed before the hidden variants of exam $k$ are revealed; it may depend on the outcomes of earlier exams. Let $C=\mathbb E[\tau_J]$ be the expected number of variants spent on the helpful version's exam.
\emph{(a)} Suppose a gate's exam $k$ depends only on $k$, the gate controls lineage error at $\alpha$, and every exam has power at least $1-\beta\ge1/2$ against $P$. Let $\bar\alpha q_k$ be the probability that exam $k$ admits a version with i.i.d.\ scores from the least-favourable law $Q^*$ of Proposition~\ref{thm:onexam}(c), so that $\sum_kq_k\le1$ by Proposition~\ref{thm:lineage}(b). Then
\[
C\ \ge\ \frac{(1-\beta)\{\log(1/\bar\alpha)+H(\pi)+\mathrm{KL}(\pi\,\|\,q)\}-h(\beta)}{K_{\inf}(P,m_0)},
\]
where $h(\beta)=-\beta\log\beta-(1-\beta)\log(1-\beta)$ and $\mathrm{KL}(\pi\|q)=\sum_k\pi_k\log(\pi_k/q_k)\ge0$.
\emph{(b)} Spending $a_k=\alpha q_k$ with $\sum_kq_k\le1$ and the bet of Proposition~\ref{thm:onexam}(b) gives a lineage-valid gate with
\[
C\ \le\ \frac{\log(1/\alpha)+H(\pi)+\mathrm{KL}(\pi\,\|\,q)+\log(1/m_0)}{g_P(\lambda)}.
\]
\end{theorem}
Compared with an exam that ignores resubmission, the honest author pays $H(\pi)/K_{\inf}$ extra variants, and among lineage-valid gates whose exams are fixed in advance and have power one, none can charge less, up to the additive constants of the theorem. With power $1-\beta<1$ the lower bound loses the factor $1-\beta$ and the term $h(\beta)$, so the gap between the two parts grows by $\beta\{\log(1/\bar\alpha)+H(\pi)+\mathrm{KL}(\pi\|q)\}+h(\beta)$, which grows with $H(\pi)$. The two parts meet for the gate of part (b) with the Kelly bet $\lambda^*=\arg\max_\lambda g_P(\lambda)$, including the boundary case $\lambda^*=1/m_0$, which occurs for \NThAllIn\% of the helpful SkillsBench skills (E11). There a variant can shrink the wealth by an arbitrary factor, but the overshoot at the crossing is still at most $\log(1/m_0)$, so Proposition~\ref{thm:onexam}(b) holds at $\lambda=1/m_0$ (Appendix~\ref{app:theory}). Its exams have power one, so $\beta=h(\beta)=0$, and exam $k$ admits $Q^*$ with probability at most $\alpha q_k$. Hence
\begin{equation}\label{eq:sandwich}
\log\frac1\alpha+H(\pi)+\mathrm{KL}(\pi\,\|\,q)\ \le\ K_{\inf}\,C\ \le\ \log\frac1\alpha+H(\pi)+\mathrm{KL}(\pi\,\|\,q)+\log\frac1{m_0}.
\end{equation}
The schedule therefore changes the honest author's cost only through $\mathrm{KL}(\pi\|q)/K_{\inf}$, to within $\log(1/m_0)/K_{\inf}$ variants that depend on neither $q$ nor $\pi$. Reading a spending sequence as a prior is old: Occam-style bounds spend $\delta P(c)$ on hypothesis $c$ with $P$ called a prior \citep[Theorem~4.1]{langford2005tutorial}, the spending functions of \citet{howard2021timeuniform} correspond to mixture densities, and weighted multiple testing uses prior weights \citep{genovese2006weighting,roeder2009genome}. Optimal alpha spending for the expected time to signal has also been derived for group-sequential binomial designs \citep{silva2020optimal}. To our knowledge, the characterization of the honest author's expected sample cost by the entropy of the attempt-count law, with a matching lower bound over lineage-valid gates whose exams are fixed in advance, is new.

\paragraph{What the theorem covers.} Three restrictions matter. First, the lower bound concerns gates whose exam $k$ is fixed before the lineage starts. Gates that carry evidence or levels from one version to the next, or across lineages through e-LOND, are outside it (Appendix~\ref{app:theory}, Scope). Second, the author may react to earlier exams, for example by submitting the helpful version sooner after a near miss, but cannot see an exam's variants before choosing what to submit to it; commit-reveal enforces exactly this. The bound then holds with $\pi$ the marginal law of $J$. The condition cannot be dropped: in simulation (E16, \NAdSkills{} SkillsBench skills) both bounds hold for authors who react to past exams, with or without the $N_{\max}$ cap, while an author who could see the variants of the first three exams and submit to the fastest one spends up to \NAdClairBelowMax\% fewer variants than the lower bound allows. Third, $C$ counts only the helpful version's exam. The useless versions also spend variants, but each of their exams is stopped by the futility bet at its own fixed level $\alpha_{\mathrm f}$ (Section~\ref{sec:gate}), so its cost does not depend on the spending schedule except through an admission probability of at most $\alpha q_k$. The entropy term is therefore the part of the cost that the schedule controls. The gate of part (b) itself has neither cap nor futility bet, so under a useless version its exam may never end; the deployed gate adds both. Futility also runs in the helpful exam: it fires with probability at most $\alpha_{\mathrm f}$ when $\Delta\ge\delta_{\mathrm f}$, but with high probability when $\delta<\Delta<\delta_{\mathrm f}$, which is one more source of $\beta>0$. We take exam $J$ to run whatever happened before; if a useless version is admitted first, which has probability at most $\alpha$, $C$ is defined on the event that exam $J$ runs. With a cap, $C$ counts only the admitting exam, not missed helpful exams that ran to $N_{\max}$. Finally, if $J$ reacts to earlier outcomes its law may depend on the gate's levels; comparisons between gates at the same $\pi$ assume it does not, as when the author reacts only to futility outcomes. In live studies 2 and 3 with Haiku each placebo and the leak skill used 186 to 209 variants before futility, against 28 and 52 for the helpful skills, so the schedule governs a small share of a lineage's total cost when most versions are useless. Two further limits concern the gate we evaluate. Power at least $1-\beta$ in every exam is a property of the untruncated exam; a cap $N_{\max}$ lowers power, and the bound applies to exams that reach the stated power before the cap. And the deployed level is $\min(\alpha_t,\alpha\gamma^{\text{lin}}_k)$: over long histories the e-LOND term sets it in \NEbElondBinds\% of first attempts (E3b), and it set every level in the live resubmission study. The theorem prices the lineage term, which binds in short histories, for early lineages and after many admissions. E19 measures the whole cost of a lineage under the deployed rule: once a lineage starts 50 or more exams into the registry, the e-LOND term sets the level in all but at most \NLcBindMax\% of exams, the choice of lineage schedule changes the helpful exam's cost by under \NLcSchedSpread\%, and the useless versions cost more than the helpful one (Appendix~\ref{app:more}). The entropy price is therefore the price of the lineage guarantee taken alone. In a long-running registry with online FDR control, resubmission protection is nearly free, and the $\log(1/\text{level})$ form of the theorem, applied to the deployed level, predicts the cost of a helpful exam to within \NLcPredErrMean\% on average; the price itself is what a new registry, or one that controls each lineage separately, pays. The simulations and reanalyses in Section~\ref{sec:results} use the grid bet, which the theory assumes and which we recommend for deployment (E13); the PrPl bet with cap $1/2$, common in earlier work, is reported beside it. The grid bet is faster at moderate and large lifts but has less power at small ones (Appendix~\ref{app:more}). Two parts still use PrPl: live studies 1 to 3, which committed it before we adopted the grid bet (studies 4 and 5 committed the grid bet), and the estimation check (E4), whose confidence bound is built from PrPl bets.

\begin{corollary}[No schedule is cheap for every attempt]\label{cor:universal}
For any $q$ with $\sum_kq_k\le1$, $\log(1/q_k)>\log k+\log\log k$ for infinitely many $k$. The schedule $q_k=c_u/\{(k+1)\log^2(k+1)\}$, with $c_u=1/\sum_{k\ge1}\{(k+1)\log^2(k+1)\}^{-1}=\NCuValue$, pays $\log k+2\log\log k+O(1)$.
\end{corollary}
The corollary is classical \citep{elias1975universal,rissanen1983universal}, and \citet{foster2013risk} already use this universal shape to set alpha-investing rules; we include it to translate it into variants.
When nothing is known about $\pi$, some attempt index always costs about $(\log k+\log\log k)/K_{\inf}$ extra variants. The schedule used in this paper, $q_k=k^{-1.6}/\zeta(1.6)$, has the shape of the spending sequences of \citet{howard2021timeuniform} and pays $1.6\log k+0.83$ nats. That is more than the universal schedule for large $k$, but less for every $k$ from 2 to 20; only the first attempt is cheaper under the universal schedule (0.71 against 0.83 nats).

\paragraph{Choosing the schedule.} By Eq.~\eqref{eq:sandwich}, a registry that does not know $\pi$ faces a coding problem: each lineage pays $\mathrm{KL}(\pi\|q)/K_{\inf}$ variants above the entropy floor. The registry may know only a class $\Pi$ that contains $\pi$, or it may learn $\pi$ from lineages that have already closed. For a class $\Pi$ of laws on the attempt indices, let
\[
R(\Pi)=\inf_q\ \sup_{\pi\in\Pi}\ \mathrm{KL}(\pi\,\|\,q),
\]
with the infimum over probability laws $q$. Write $\Pi_N$ for all laws on $\{1,\dots,N\}$, $\Pi_N^{\downarrow}$ for the non-increasing ones, and $C(\pi)$ for the cost of Theorem~\ref{thm:price} when $J\sim\pi$.

\begin{proposition}[Minimax schedule]\label{prop:minimax}
\emph{(a)} Every gate as in Theorem~\ref{thm:price}(a) has
\[\sup_{\pi\in\Pi}\big\{K_{\inf}C(\pi)-(1-\beta)[\log(1/\bar\alpha)+H(\pi)]\big\}\ \ge\ (1-\beta)R(\Pi)-h(\beta).\]
\emph{(b)} If $\sup_{\pi\in\Pi}\mathrm{KL}(\pi\|q)\le R(\Pi)+\epsilon$, the gate of Theorem~\ref{thm:price}(b) with the Kelly bet $\lambda^*$ has
\[\sup_{\pi\in\Pi}\big\{K_{\inf}C(\pi)-\log(1/\alpha)-H(\pi)\big\}\ \le\ R(\Pi)+\epsilon+\log(1/m_0).\]
\emph{(c)} For any $\pi_1,\dots,\pi_M\in\Pi$ and weights $w$, $R(\Pi)\ge\sum_jw_j\,\mathrm{KL}(\pi_j\|\bar\pi_w)$, where $\bar\pi_w=\sum_jw_j\pi_j$.
\emph{(d)} $R(\Pi_N)=\log N$, attained by the uniform schedule, and $R(\Pi_N^{\downarrow})\le\log H_N$ with $H_N=\sum_{k\le N}1/k$, the worst-case excess of $q_k=1/(kH_N)$.
\end{proposition}
Part (c) is the easy half of the redundancy-capacity theorem, by which $R(\Pi)$ equals the capacity $\sup_w\sum_jw_j\mathrm{KL}(\pi_j\|\bar\pi_w)$ of the channel $j\mapsto\pi_j$ under general conditions, with the capacity-achieving mixture as the minimax schedule \citep{haussler1997general,merhav1995strong}. Parts (a) and (b) say that this capacity, divided by $K_{\inf}$, is the price in variants of knowing only the class of $\pi$. Over all laws on the positive integers $R=\infty$, which is Corollary~\ref{cor:universal} again; a finite price needs a restricted class, as in coding on countable alphabets \citep{boucheron2009coding}. For natural classes the price is small (Appendix~\ref{app:more}, E12). If honest authors make at most 20 attempts, the minimax excess is $\log20=\NScLogTwenty$ nats, against \NScZetaWorstTwenty{} nats in the worst case for the deployed $\zeta(1.6)$ schedule. If, in addition, later attempts are never more likely than earlier ones, the excess for $N=1{,}000$ lies in the certified interval $[\NScMonoLoThousand,\NScMonoHiThousand]$ nats, computed with the Blahut-Arimoto algorithm \citep{blahut1972computation,arimoto1972algorithm} on the extreme points of $\Pi_N^{\downarrow}$. That is against $\log H_N=\NScLogHThousand$ for the explicit schedule and $\log N=\NScLogThousand$ without the monotonicity.

\begin{theorem}[Learned schedule]\label{thm:learned}
Close every lineage after at most $N$ attempts. When lineage $i$ opens, the registry fixes
\[
\hat q^{(i)}_k=(1-\epsilon)\,\frac{c_k+1}{m+N}+\frac{\epsilon}{N},\qquad k=1,\dots,N,
\]
where $c_k$ is the number of counted lineages admitted at attempt $k$, and $m=\sum_kc_k$. The counted lineages must be closed when lineage $i$ opens and must be chosen without regard to their attempt indices. Counting every lineage that happens to have closed is not enough, because lineages that need fewer attempts close sooner and would be over-represented. The rule used below counts lineages $1,\dots,i-1-D$, which is valid when every lineage closes before $D$ further lineages open. Exam $k$ of lineage $i$ runs at level $\alpha\hat q^{(i)}_k$.
\emph{(a)} If lineage $i$ is null in the sense of Definition~\ref{def:condnull}, it is ever admitted with probability at most $\alpha$, whatever the other lineages do, including lineages run by the same author to move the counts.
\emph{(b)} Suppose that $m_h$ of the counted lineages are honest, with attempt indices i.i.d.\ $\pi\in\Pi_N$, and that at most $r$ other counts are arbitrary and may depend on them. Then
\[
\mathbb E\,\mathrm{KL}(\pi\,\|\,\hat q^{(i)})\ \le\ \log\frac1{1-\epsilon}+\log\Big(1+\frac{N-1+r}{m_h+1}\Big),
\]
and $\mathrm{KL}(\pi\|\hat q^{(i)})\le\log(N/\epsilon)-H(\pi)$ for every value of the counts.
\emph{(c)} Suppose lineages open one after another, all are honest and recorded at their attempt index, every lineage closes before $D$ further lineages open, and lineage $i$ counts lineages $1,\dots,i-1-D$. Then $\sum_{i\le n}\mathbb E\,\mathrm{KL}(\pi\|\hat q^{(i)})\le n\log\frac1{1-\epsilon}+(D+1)\log N+(N-1)\log n$. Conversely, for any rule that fixes the schedule of lineage $i$ from the attempt indices of lineages $1,\dots,i-1$,
\[\sup_{\pi\in\Pi_N}\sum_{i\le n}\mathbb E\,\mathrm{KL}(\pi\,\|\,q^{(i)})\ \ge\ R_n(N)=\frac{N-1}2\log\frac n{2\pi e}+\log\frac{\Gamma(1/2)^N}{\Gamma(N/2)}+o(1),\]
the minimax redundancy of $n$ i.i.d.\ draws from $\Pi_N$ \citep{xie1997minimax}.
\end{theorem}
\emph{Capped exams.} With a cap $N_{\max}$, a helpful version can miss admission; the lineage is then admitted at a later attempt or closes unadmitted. We define $c_k$ as the number of counted lineages \emph{admitted} at attempt $k$, so $\hat q$ learns $\tilde\pi$, the law of the admitting attempt among admitted lineages. This is also the law that matters for cost, because Theorem~\ref{thm:price} charges the exam that admits. Part (a) does not use the counts and is unchanged, and parts (b) and (c) hold with $\tilde\pi$ in place of $\pi$ (Appendix~\ref{app:theory}). If each helpful exam at attempts $k\le N$ admits before the cap with probability at least $1-\beta_{\max}$, then, with $\varepsilon=\beta_{\max}+\alpha$, $\tilde\pi$ and $\pi$ differ in total variation by at most $\varepsilon+\varepsilon/(1-\varepsilon)$. A uniform $\beta_{\max}$ is only plausible because lineages close after $N$ attempts: wealth after $N_{\max}$ variants is at most $m_0^{-N_{\max}}$, so no exam admits at a level below $m_0^{N_{\max}}$, and with unlimited attempts $q_k\to0$ would make late helpful exams powerless.

Learning the schedule costs nothing in validity. A null author who floods the registry with lineages of their own can raise the honest authors' cost, by at most $\log(N/\epsilon)$ nats per lineage, but cannot raise any lineage's error. Learning levels from earlier results while keeping validity is also the idea of contextual online FDR control \citep{chen2020contextual}; what Theorem~\ref{thm:learned} adds is the price. The excess per lineage falls like $(N-1)/m$, and with $\epsilon=1/n$ the cumulative excess is, as $n\to\infty$ with $N$ fixed, within a factor of two of the minimax redundancy, so no rule learns much faster. The factor is asymptotic: at $N=20$ and $n=400$ the bound is about four times $R_n(N)$. It comes from the add-one estimator; an add-one-half (Krichevsky-Trofimov) estimator would match the leading term. The cost converse assumes the same $P$, and so the same $K_{\inf}$, in every lineage. Lineages that overlap in time cost only the additive $(D+1)\log N$, provided the registry waits for a fixed number of later openings rather than counting whichever lineages have closed. In E12b, with one attempt per time step and $N=20$, the naive rule over-represents short lineages: with uniform $\pi$ it puts mass \NObUnifQoneNaive{} on the first attempt after 25 lineages, against \NObUnifQoneDelay{} for the fixed-delay rule and $\pi_1=\NObUnifPiOne$. Its cumulative excess after 400 lineages is higher when long lineages are common (\NObIncNaive{} against \NObIncDelay{} nats for $\pi_k\propto k$) and lower when short ones dominate (\NObGeoNaive{} against \NObGeoDelay{} for the geometric law), because it uses more lineages. Only the fixed-delay rule carries the guarantee, and it stays far below the bound of part (c) (\NObBound{} nats). With power-one exams, Theorem~\ref{thm:price}(a) applies to each lineage given the earlier ones with $\beta=0$, so the converse holds for the cost itself: some $\pi\in\Pi_N$ gives $\sum_{i\le n}\{K_{\inf}\mathbb E[C_i]-\log(1/\bar\alpha)-H(\pi)\}\ge R_n(N)$. In E12, with $N=20$, $\epsilon=0.05$ and geometric attempt counts with mean 3, the expected excess falls from \NScKLZero{} nats with no closed lineages to \NScKLHundred{} after 100 and \NScKLThousand{} after 1{,}000, below the bound at every $m$. With 50 poisoned counts placed where $\pi$ is smallest, it is \NScKLPoisonHundred{} nats after 100 honest lineages, against the bound \NScBoundPoisonHundred. For the median skill, simulated extra variants match $\mathrm{KL}(\pi\|\hat q)/g_P(\lambda^*)$: \NScVarHundredSim{} against \NScVarHundredTh{} after 100 closed lineages and \NScVarFiveHundredSim{} against \NScVarFiveHundredTh{} after 500. The deployed $\zeta(1.6)$ schedule, renormalized to 20 attempts, costs \NScVarZetaSim{} extra variants for this $\pi$, so the learned schedule overtakes it between 100 and 500 closed lineages.

\paragraph{Not knowing the skill.} Theorem~\ref{thm:price}(b) bets with a $\lambda$ chosen for $P$, which a registry does not know. A mixture of fixed bets removes this need, and spending makes a natural choice of grid available: each exam has its own level, fixed in advance, so the grid can be tuned to it. For an exam at level $a$, let $L=\log(1/a)$, $\rho=L/(L+1)$ and $\lambda_j=\rho^{j+1}/m_0$ for $j=0,1,\dots$, with the weights $w_j=c_u/\{(j+2)\log^2(j+2)\}$ of Corollary~\ref{cor:universal}, which sum to one. The exam admits at the first $n$ with
\[
M_n=\sum_{j\ge0}w_j\prod_{i\le n}\{1+\lambda_j(X_i-m_0)\}\ \ge\ 1/a .
\]
Let $\lambda^*\in(0,1/m_0]$ maximize $g_P$ on $[0,1/m_0]$, so that $g_P(\lambda^*)=K_{\inf}(P,m_0)$, and write $\Lambda=\log\{1/(m_0\lambda^*)\}\ge0$ for how far the Kelly bet lies below the largest admissible bet.
\begin{theorem}[Level-tuned grid]\label{thm:adaptive}
\emph{(a)} The exam does not depend on $P$, is valid at level $a$ under Definition~\ref{def:condnull}, and has power one against every $P$ with $\mathbb E_PX>m_0$.
\emph{(b)} For every such $P$, with $j^*=\lfloor\Lambda/\log(1/\rho)\rfloor$ and $\ell(j)=\log(1/w_j)$,
\[
K_{\inf}(P,m_0)\,\mathbb E_P[\tau]\ \le\ \Bigl(1+\frac1L\Bigr)\Bigl\{L+\log\frac1{m_0}+\ell(j^*)\Bigr\}.
\]
With $u=\log\{(L+1)(2+\Lambda)\}$, this is at most $L+1+(1+1/L)\{\log(L+1)+\log(2+\Lambda)+2\log u+\log(1/m_0)+\log(1/c_u)\}$.
\emph{(c)} Spending $a_k=\alpha q_k$ with this exam at each level gives a lineage-valid gate that does not depend on $P$. With $\bar L=\log(1/\alpha)+H(\pi)+\mathrm{KL}(\pi\|q)$ and $\bar u=\log\{(\bar L+1)(2+\Lambda)\}$,
\[
K_{\inf}C\ \le\ \bar L+1+\Bigl(1+\frac1{\log(1/\alpha)}\Bigr)\Bigl\{\log(\bar L+1)+\log(2+\Lambda)+2\log\bar u+\log\frac1{m_0}+\log\frac1{c_u}\Bigr\},
\]
while every lineage-valid gate whose exam $k$ depends only on $k$ and has power one has $K_{\inf}C\ge\log(1/\bar\alpha)+H(\pi)+\mathrm{KL}(\pi\|q')$, with $q'$ its own size schedule under $Q^*$.
\end{theorem}
In part (c) we replaced $1+1/\bar L$ by $1+1/\log(1/\alpha)$, which is valid because the braced term is nonnegative: $m_0\le1$, $c_u<1$, and each exam's $u_k=\log\{(L_k+1)(2+\Lambda)\}\ge1$ because $L_k\ge\log(1/\alpha)\ge\log2$ for every $\alpha\le0.5$, which covers every level used here. In part (c) the lower bound uses the gate's own schedule $q'$ and the upper bound uses $q$. For the grid gate, exam $k$ admits $Q^*$ with probability at most $\alpha q_k$, so $\bar\alpha q'_k\le\alpha q_k\le\bar\alpha q_k$ and $\mathrm{KL}(\pi\|q')\ge\mathrm{KL}(\pi\|q)$; the two bounds therefore share the leading terms $\log(1/\alpha)+H(\pi)+\mathrm{KL}(\pi\|q)$ up to $\log(\bar\alpha/\alpha)$. So the entropy price survives without knowing the skill: at every level the leading term of $K_{\inf}\mathbb E[\tau]$ is $L$ with coefficient one, and not knowing $P$ costs at most about $\log\bar L+\log\Lambda$ nats. The construction combines known parts, and its first-order behaviour is known. Mixing over a geometric grid whose resolution is tuned to the level is the device of \citet[proof of Theorem~33]{kaufmann2021mixture}, who apply it to grids over time in deviation inequalities; here we move it to the bet. For bounded payoffs, \citet[Theorem~2]{orabona2026optimal} characterizes the expected rejection time of portfolio tests exactly as $\alpha\to0$: with a prior whose log-scale density decays at the integrability border, which needs no knowledge of the level, $\mathbb E_P[\tau]$ is asymptotic to $\{\log(1/\alpha)+\log\log(1/|\lambda^*_P|)\}/K_{\inf}$, with $\lambda^*_P$ the Kelly bet under $P$, uniformly over alternatives, and his Theorem~5 shows that the iterated-logarithm term cannot be removed. Asymptotically his bound is sharper than ours, which pays an extra $\log L$. What Theorem~\ref{thm:adaptive} adds is a bound that holds at every finite level with explicit constants, which a registry needs for exams at levels such as $\alpha q_k$ that are small but fixed, and part (c), which combines it with the spending schedule; to our knowledge that combination is new. For exponential families and in an integrated (Bayes) sense, \citet{pollak1978optimality} shows that the second-order term in $L$ cannot be smaller than $\tfrac12\log L$, so our coefficient of about one on $\log L$ is not optimal. The term in $\Lambda$ has the order of the law of the iterated logarithm, which cannot be avoided when the effect is small \citep{farrell1964asymptotic,agrawal2025stopping,orabona2026optimal}.
In E13 (Appendix~\ref{app:more}) we ran the grid exam on the \NThHelpful{} SkillsBench skills with positive released lift. At the first exam's level $\alpha q_1$ it needs a median of \NAgGridQone{} variants, against \NAgKellyQone{} for the Kelly bet that knows $P$, \NAgAgrapaQone{} for aGRAPA and \NAgPrplQone{} for the PrPl bet with cap $1/2$; the bound of part (b) gives \NAgBoundQone. For the median skill, $K_{\inf}\mathbb E[\tau]-L$ stays between \NAgExcessMin{} and \NAgExcessMax{} nats as $L$ grows from $\log(1/\alpha)$ to $\log(1/\alpha)+16$, against \NAgFixedMin{} to \NAgFixedMax{} nats for the same grid with $\rho=1/2$ fixed. Under each skill's least-favourable null, the largest rejection rate is \NAgMaxReject\% at level \NAgLevelQone\%.

\paragraph{What the results say about the gate.} We checked each statement on the \NThHelpful{} SkillsBench skills with positive released lift, taking each skill's 18 per-configuration differences as $P$ (Appendix~\ref{app:more}, E11).
\begin{itemize}
\item \emph{Duality.} The least-favourable law $Q^*$ of Proposition~\ref{thm:onexam}(c) is a null law with $\mathrm{KL}(P\|Q^*)=K_{\inf}$ for every skill, to numerical precision.
\item \emph{One exam.} No simulated cost falls below the lower bound. Relative to $\log(1/\alpha)/K_{\inf}$, the Kelly bet needs a median $\NThOracleRatio\times$ (at most $\NThOracleRatioMax\times$), aGRAPA $\NThAgrapaRatio\times$, and the PrPl bet with $c=1/2$ (used in live studies 1 to 3) needs $\NThPrplRatio\times$ (at most $\NThPrplRatioMax\times$). For \NThAllIn\% of these skills the Kelly bet is $\lambda^*=1/m_0$, and the cap is what costs the most variants. This is why the deployed gate uses the level-tuned grid bet.
\item \emph{Lineage.} For the median skill ($K_{\inf}=\NThMedKinf$, about \NThMedIdeal{} variants at level $\alpha$), the simulated extra cost of the $k$-th attempt under our schedule matches $\log(q_1/q_k)/g_P(\lambda^*)$: \NThExtraTwoSim{} against \NThExtraTwoTh{} variants at $k=2$, \NThExtraTenSim{} against \NThExtraTenTh{} at $k=10$, and \NThExtraTwentySim{} against \NThExtraTwentyTh{} at $k=20$.
\item \emph{Honest authors.} An author who always succeeds on the first attempt pays \NThPriceOne{} extra variants, because the first exam runs at $\alpha q_1=0.022$. If honest attempt counts are geometric with mean 3, the expected price is \NThPriceGeoThree{} variants against the entropy floor $H(\pi)/K_{\inf}=\NThFloorGeoThree$. If they follow a Zipf law with exponent 2, it is \NThPriceZipf{} against \NThFloorZipf{}. A geometric schedule fits geometric authors better, but costs \NThGeoOnZipfLoss{} more variants under the Zipf law.
\end{itemize}
